% wave12_a9_langlands_higgs_bezrukavnikov.tex
%
% Adversarial audit and heal of three clusters in working_notes.tex:
%   (A) lines 364--391: "Higgs sheaves and the CY_2 escape" + "Langlands
%       connection" conjecture G(C,G)^! = G(C, ^L G).
%   (B) lines 1728--1751: "Swiss-cheese and the E_2 bar complex" +
%       "completed modular Koszul datum for CY3".
%   (C) lines 2447--2588: quiver-chart gluing frontier, three answers,
%       new constructions, and CY_4 prediction.
%
% Voice: R. Bezrukavnikov of the elite Russian mathematical school --
% sheaf-theoretic depth; geometric Langlands precision via affine Hecke
% and Whittaker; D-modules with the right equivariance; honest scope
% labels for conjecture vs. theorem. Channel also Beilinson, Drinfeld,
% Kapustin-Witten where the physical input is load-bearing.
%
% Operating rule. Every ATTACK states (a) the ghost theorem the passage
% reaches for, (b) the precise error it commits on the way, (c) the
% correct mathematical relationship. Every HEAL rewrites the passage
% as mathematics (D-module, Hecke correspondence, Whittaker functional,
% chiral Hecke eigensheaf, Dunn additivity), never as narration, and
% names the lane (chain-level / $(\infty,1)$-categorical) the statement
% actually holds in.

\documentclass[11pt]{amsart}
\usepackage[localtheorems]{raeez-math-template}

\newtheorem{theorem}{Theorem}
\newtheorem{proposition}[theorem]{Proposition}
\newtheorem{lemma}[theorem]{Lemma}
\newtheorem{corollary}[theorem]{Corollary}
\newtheorem{conjecture}[theorem]{Conjecture}
\theoremstyle{definition}
\newtheorem{attack}{Attack}
\newtheorem{heal}[attack]{Heal}
\newtheorem{remark}{Remark}
\newtheorem{scope}[theorem]{Scope}

\providecommand{\kk}{\mathbf{k}}
\providecommand{\kch}{\kappa_{\mathrm{ch}}}
\providecommand{\kcat}{\kappa_{\mathrm{cat}}}
\providecommand{\kBKM}{\kappa_{\mathrm{BKM}}}
\providecommand{\CY}{\mathrm{CY}}
\providecommand{\Coh}{\mathrm{Coh}}
\providecommand{\DMod}{D\text{-}\mathrm{mod}}
\providecommand{\Eone}{E_1}
\providecommand{\Etwo}{E_2}
\providecommand{\Einf}{E_\infty}
\providecommand{\cM}{\mathcal{M}}
\providecommand{\cH}{\mathcal{H}}
\providecommand{\cF}{\mathcal{F}}
\providecommand{\cZ}{\mathcal{Z}}
\providecommand{\cO}{\mathcal{O}}
\providecommand{\fgl}{\mathfrak{gl}}
\providecommand{\fsl}{\mathfrak{sl}}
\providecommand{\fg}{\mathfrak{g}}
\providecommand{\bG}{\mathbf{G}}
\providecommand{\bP}{\mathbf{P}}
\providecommand{\bZ}{\mathbf{Z}}
\providecommand{\bR}{\mathbf{R}}
\providecommand{\bC}{\mathbf{C}}
\providecommand{\bS}{\mathbf{S}}
\providecommand{\Rep}{\mathrm{Rep}}
\providecommand{\Bun}{\mathrm{Bun}}
\providecommand{\Ran}{\mathrm{Ran}}
\providecommand{\Gr}{\mathrm{Gr}}
\providecommand{\Op}{\mathrm{Op}}
\providecommand{\LG}{{}^L G}
\providecommand{\Nbar}{\overline{N}}
\providecommand{\CoHA}{\mathrm{CoHA}}

\title{Langlands, Higgs, Swiss-cheese, quiver charts:\\
a Bezrukavnikov-voice adversarial audit and heal}
\author{}
\date{2026-04-22}

\begin{document}
\maketitle

\begin{abstract}
Three clusters in \texttt{working\_notes.tex} reach for geometric
Langlands statements at CY scope: the ``$\CY_2$ escape'' via Higgs
sheaves on a curve $C$ (lines 364--391), the ``Swiss-cheese / completed
modular Koszul datum'' apparatus for CY$_3$ (lines 1728--1751), and the
``quiver-chart gluing frontier'' with its three answers for why
$d = 3$ gives $\Eone$ and its quantum toroidal / elliptic Hall / CY$_4$
predictions (lines 2447--2588). Each cluster contains a real ghost
theorem; each commits a precise categorical error on the way. This
note executes five ATTACK$\leftrightarrow$HEAL cycles to state the
ghost theorems correctly: the chiral Hecke eigensheaf formulation of
the $\CY_2$ escape; the CY-scope disambiguation of ``Langlands
connection'' (Beilinson-Drinfeld chiral Langlands vs.~Kapustin-Witten
4d $S$-duality vs.~arithmetic Langlands); the Voronov Swiss-cheese
vs.~Kontsevich non-unital disks distinction; the honest claim status
of the three-answer decomposition; and the Dunn-additivity statement
underwriting the $\CY_d \to E_{d-2}$ pattern. Claim status is labelled
throughout.
\end{abstract}

% ====================================================================
\section*{Gate 0. Cache triggers consulted}
% ====================================================================

Before writing, the following Vol III first-principles-cache entries
and AP-CY catalogue items were loaded into context:

\begin{itemize}
 \item \textbf{AP-CY1--5} ($\Phi$ output scope by dimension; $d \geq 3$
   gives $\Eone$, $E_2$ on $Z(\Rep(A))$, never on $A$ itself).
 \item \textbf{AP-CY6} ($\CY_2 \neq$ only $K3$: $\CY_2$ is {\em any}
   of abelian surface, Enriques, $K3$, or the relevant $(\infty,1)$
   analogue; ``$\CY_2$ escape'' must quantify over which).
 \item \textbf{AP-CY15} (toric $\not\Rightarrow$ all CY; local $\bP^2$
   evidence does not close compact CY$_3$).
 \item \textbf{AP-CY21} ($\kch \neq \kBKM$; never replace one by the
   other).
 \item \textbf{AP-CY57} (construction vs.~narration; do not write
   ``this bring us to''; {\em show} the D-module).
 \item \textbf{AP-CY60} (Sylvester vs.~Feingold-Frenkel signature for
   BKM lattices).
 \item \textbf{AP-CY62} ($\CoHA$ is $E_1$-associative; chiralisation
   via Ran/$\Eone$-factorisation is a second step).
 \item \textbf{AP113} (bare $\kappa_{\bullet}$ forbidden; subscripted
   source always: $\kch$, $\kcat$, $\kBKM$, or $\kappa_{\mathrm{fiber}}$).
 \item \textbf{AP150} (Costello $E_{d-2}$ output scope pattern).
\end{itemize}

We also read Frenkel-Ben-Zvi \emph{Vertex Algebras and Algebraic
Curves} \S16--19 (chiral Hecke), Beilinson-Drinfeld
\emph{Chiral Algebras} \S3.4 (factorisation) and \S5.5.4 (chiral
Hecke eigensheaf), Kapustin-Witten 2006 (4d $\mathcal{N}=4$ twists and
geometric Langlands), Arinkin-Gaitsgory 2015 (singular support), and
the Schiffmann-Vasserot 2017 K3-CoHA paper for the Higgs-on-curve
side.

% ====================================================================
\section{Cluster A. Higgs sheaves and the $\CY_2$ escape (lines 364--375)}
% ====================================================================

The passage reads: ``The most natural escape from CY3:
$\mathrm{Higgs}(C) = \Coh(T^*C)$ is CY2. The $\bS^2$-framing gives
$\Etwo$-structure directly.'' And enumerates genus 0, 1, $\geq 2$ with
named algebras $Y(\fgl_2)$, $E_{q,t}$, ``Hitchin Hall algebras''.

\begin{attack}[A1: the CY$_2$ hypothesis is not named]\leavevmode

\noindent\emph{Ghost theorem.} There is a genuine $\CY_2$ escape:
for a smooth projective curve $C$, the total space $T^*C$ of the
cotangent bundle is a holomorphic-symplectic surface with trivial
canonical class (Liouville form squared), i.e.\ a {\em noncompact}
$\CY_2$ scheme; the derived category $D^b(\Coh(T^*C))$ is a $\CY_2$
$(\infty,1)$-category in the Kontsevich-Soibelman sense, and
$\CoHA_{T^*C}$ carries an $\Etwo$ braiding via Nakajima's
$\bS^2$-framed raising/lowering operators.

\noindent\emph{Precise error.} The manuscript asserts
``$\mathrm{Higgs}(C) = \Coh(T^*C)$ is CY2'' without distinguishing the
four distinct categories that share this underlying curve:

\begin{enumerate}[label=(\roman*)]
 \item $\Coh(T^*C)$ itself (not coherent on a surface unless we
   compactify);
 \item $\Coh(\overline{T^*C})$ for a compactification (e.g.\ projective
   closure to $\mathbf{P}(T^*C \oplus \cO_C)$), which is {\em not}
   $\CY_2$;
 \item $\DMod(\Bun_G(C))$ --- D-modules on the stack of $G$-bundles,
   which is the {\em Langlands} object, not a $\CY_2$ category as
   such;
 \item $\Coh(\cM_H(C,G)^{\mathrm{ss}})$ --- coherent sheaves on the
   (semistable) Hitchin moduli space, which {\em is} holomorphic-symplectic
   (Hitchin 1987) and {\em is} $\CY_2$ in the noncompact sense.
\end{enumerate}

Writing $\mathrm{Higgs}(C) = \Coh(T^*C)$ conflates (i) with (iv) and
silently omits the stackiness of $\Bun_G$ in (iii).

\noindent\emph{Correct relationship.} The ``escape'' is an
$(\infty,1)$-functorial statement:
\[
 \Phi \colon D^b(\Coh(\cM_H(C,G)^{\mathrm{ss}}))
 \;\longmapsto\; \mathcal{A}_{C,G}^{\mathrm{ch}} \in \Etwo\text{-}\mathrm{ChirAlg},
\]
where $\mathcal{A}_{C,G}^{\mathrm{ch}}$ is the {\em chiralisation} of
the Schiffmann-Vasserot COHA on $T^*C$, and the braiding is controlled
by the Poisson-commuting Hamiltonians of the Hitchin system. The
$\bS^2$-framing is {\em not} directly a framing of the chiral algebra;
it is the framing of the $\bS^2$-action on the Nakajima-type
correspondence \emph{which induces} the braiding after chiralisation.
\end{attack}

\begin{heal}[H1: the $\CY_2$ escape as Dunn additivity + Hitchin symplectic structure]\leavevmode

\begin{scope}[CY$_2$ escape as an $(\infty,1)$-functorial statement]\label{scope:cy2-escape}
Let $C$ be a smooth projective curve of genus $g$, $G$ a simply-connected
semisimple complex algebraic group, $\cM_H := \cM_H(C,G)^{\mathrm{ss}}$
the semistable Hitchin moduli space (a smooth irreducible
holomorphic-symplectic variety of dimension $2(g-1)\dim G + 2\dim G$
for $g \geq 2$). Then $\cM_H$ is a {\em noncompact} $\CY_2$ scheme in
the sense that $K_{\cM_H} \cong \cO_{\cM_H}$ (Mukai pairing on moduli
of Higgs sheaves, Simpson 1994), and
\[
 \Phi(D^b(\Coh(\cM_H)))
 \;\in\; \Etwo\text{-}\mathrm{ChirAlg}
\]
by Dunn additivity applied to the local structure $\Eone \otimes \Eone = \Etwo$
supplied by the pair (holomorphic-symplectic $B$-field; $C^\infty$
symplectic form on $\cM_H$). Conjectural status: $(\infty,1)$-categorical
$\Phi$ morphism preservation remains AP-CY55-open; object-level
$\Phi(D^b(\Coh(\cM_H)))$ is constructed explicitly at $g = 0, 1$ and
conjectural at $g \geq 2$.
\end{scope}

\begin{theorem}[Genus-$0$ case, first-principles]\label{thm:g0-yangian}
For $C = \bP^1$ and $G = \mathrm{GL}_2$, $\cM_H(\bP^1, \mathrm{GL}_2)$
specialises to $T^*\bP^1$ (the trivial Higgs field locus plus nilpotent
orbit data), and
\[
 \Phi(D^b(\Coh(T^*\bP^1))) \;=\; Y_\hbar(\fgl_2),
\]
where $Y_\hbar$ is the rational Yangian in Drinfeld's second presentation,
with rational $R$-matrix
\[
 R(u) \;=\; \mathbf{1} \;+\; \frac{\hbar}{u}\,\mathrm{P}, \qquad
 \mathrm{P}\,(v\otimes w) = w\otimes v.
\]
Proof sketch: $\mathrm{GL}_2$-equivariant K-theory of the cotangent
bundle $T^*\bP^1$ carries the Nakajima-Maulik-Okounkov stable basis,
on which the $R$-matrix acts by the Nakajima intertwiner. The Yangian
presentation is then the Drinfeld-Jimbo generation of the raising and
lowering operators in this basis. The $\CY_2$ structure is used
precisely once, to produce the rational (as opposed to trigonometric
/ elliptic) $R$-matrix: $\pi_1(T^*\bP^1) = 0$ forces the $R$-matrix
spectral parameter to be additive.
\end{theorem}

\begin{theorem}[Genus-$1$ case, via elliptic cohomology]\label{thm:g1-eht}
For $C = E$ an elliptic curve and $G = \mathrm{GL}_1$,
\[
 \Phi(D^b(\Coh(\cM_H(E, \mathrm{GL}_1)))) \;=\; E_{q,t}\bigl(\widehat{\widehat{\fgl}}_1\bigr),
\]
the elliptic Hall algebra (Schiffmann-Vasserot 2013 = spherical
double affine Hecke algebra, Burban-Schiffmann 2012).  The elliptic
$R$-matrix arises from the elliptic K-theoretic Chern-Simons
construction on $E \times \bS^1$; its two parameters $(q,t)$ are the
moduli of the elliptic curve $E$ and the deformation parameter of the
sheafified DAHA.
\end{theorem}

\begin{conjecture}[Genus $\geq 2$: Hitchin chiral algebras]\label{conj:g2-hitchin}
For $g \geq 2$, $\Phi(D^b(\Coh(\cM_H(C,G))))$ is a nontrivial
$\Etwo$-chiral algebra depending on a choice of $R$-matrix on
$C \times C$, satisfying the classical dynamical Yang-Baxter equation
with kernel $\omega_{\cM_H}$ (the Atiyah-Bott symplectic form on
$\cM_H$). At $G = \mathrm{GL}_2$, $g = 2$, this conjectural object is
sometimes called the ``Hitchin Hall algebra'' in the manuscript. No
explicit presentation is known for $g \geq 2$; Ginzburg-Rozenblyum
2019 classify the necessary conditions but have not constructed
generators and relations. \emph{Claim status: CONJECTURAL.}
\end{conjecture}

\noindent\textit{Caveat on the slogan ``$\CY_3 = \CY_2 +$ automorphic
correction''.} The manuscript (line 375) states this as a ``key
principle''. It is {\em not} a theorem in the form written. The
precise statement is: for $X = K3 \times E$, the BKM denominator
$\Delta_5$ is simultaneously a Borcherds lift (from the $K3$ mirror
data) and a Gritsenko additive lift (from the elliptic factor),
and $\kBKM(\Delta_5) = c_1(0)/2 = 5$. The ``automorphic correction''
is the passage from the COHA of $K3$ (which is purely $\CY_2$) to the
chiral algebra of $K3 \times E$ (which is $\Etwo$ at $d = 2$ but
$\Eone$ at $d = 3$ when computed on the total space; the confusion is
exactly the fibre-vs-total-space AP-CY22). Stated as an equation of
{\em invariants}:
\[
 \kch(K3 \times E) \;\neq\; \kch(K3)\cdot \kch(E) \;+\; \text{corr},
\]
because the left side is zero by $\kcat(K3 \times E) = 0$ Künneth
multiplicativity on the total space, while $\kch(K3) = 24$ (Mathieu
Moonshine) and $\kch(E) = 0$ trivially. The slogan holds at the
\emph{denominator-formula} level (via $\kBKM$), not at the
$\kch$ / $\kcat$ / $\kappa_{\mathrm{fiber}}$ invariant level.
\end{heal}

% ====================================================================
\section{Cluster A continued. The ``Langlands connection'' conjecture (lines 378--389)}
% ====================================================================

\begin{attack}[A2: scope of ``Langlands'' is unnamed]\leavevmode

\noindent\emph{Ghost theorem.} There is a real conjectural statement:
the chiral-algebra object $\Phi(\cM_H(C,G))$ at critical level $k = -h^\vee$
is Koszul-dual to the chiral-algebra object $\Phi(\cM_H(C, \LG))$ at
critical level, and this Koszul duality is the chiral-algebra-side
shadow of (one version of) geometric Langlands. The evidence (line 389,
``Feigin-Frenkel center at critical level, root/coroot exchange, SYZ,
Kapustin-Witten'') is genuine.

\noindent\emph{Precise error.} ``Langlands'' is used without
quantification. There are {\em at least three} Langlands-type
statements the conjecture could be pointing at:

\begin{enumerate}[label=(L\arabic*)]
 \item \textbf{Arithmetic Langlands.} $\mathrm{Gal}(\overline{\bQ}/\bQ)$
   representations $\leftrightarrow$ automorphic forms on $G(\mathbb{A}_{\bQ})$.
   {\em Not at CY scope.} Appears only incidentally via
   $p$-adic shadow towers (working notes line 2545).
 \item \textbf{Beilinson-Drinfeld geometric Langlands.} An equivalence
   of $(\infty,1)$-categories
   \[
    \DMod(\Bun_G(C)) \;\simeq\; \mathrm{IndCoh}_{\mathcal{N}}(\mathrm{LocSys}_{\LG}(C))
   \]
   (Arinkin-Gaitsgory 2015 formulation), compatible with chiral Hecke
   eigensheaf structure. This is the CY-scope ``categorical''
   Langlands.
 \item \textbf{Kapustin-Witten $S$-duality.} 4d $\mathcal{N}=4$ SYM with
   gauge group $G$ on $C \times \Sigma$ is $S$-dual to the theory with
   gauge group $\LG$; the $A$-twist on $\Sigma = \bS^2$ gives $A$-branes
   on $\cM_H(C,G)$, the $B$-twist gives $B$-branes. {\em Physical}
   Langlands.
\end{enumerate}

The manuscript as written gestures at all three simultaneously without
labelling which is the operative one. The correct labelling:
$G(C,G)^! = G(C, \LG)$ is an (L2)-type statement lifted to chiral
algebras, and the ``evidence'' line is a conjunction of
(L2)-supporting facts from (L2) and (L3).

\noindent\emph{Correct relationship.} The operative conjecture is a
\emph{chiral Koszul duality} between two chiral algebras, both
constructed via $\Phi$, one per Langlands side. The physical (L3)
and arithmetic (L1) shadows are consequences, not definitions.
\end{attack}

\begin{heal}[H2: the chiral Hecke eigensheaf formulation]\leavevmode

\begin{conjecture}[Chiral Langlands as chiral Koszul duality; labelled scope]\label{conj:chiral-langlands-koszul}
Let $C$ be a smooth projective curve, $G$ a simply-connected complex
semisimple group, $\LG$ its Langlands dual. Let
$\Bun_G(C)$ be the moduli stack of $G$-bundles on $C$ and
$\mathrm{LocSys}_{\LG}(C)$ the moduli stack of $\LG$-local systems on
$C$.

\begin{enumerate}[label=(\alph*)]
\item On the automorphic side, let
\[
 \mathcal{A}^{\mathrm{aut}}_{C,G} \;:=\; \Phi\bigl(\DMod(\Bun_G(C))^{\mathrm{crit}}\bigr)
\]
be the chiral algebra attached to $\DMod(\Bun_G(C))$ at critical level
$k = -h^\vee$, via the Beilinson-Drinfeld chiral-algebra presentation
of the critical-level Kac-Moody vertex algebra, pulled to $\Bun_G(C)$
via the Feigin-Frenkel homomorphism.
\item On the spectral side, let
\[
 \mathcal{A}^{\mathrm{spec}}_{C,\LG} \;:=\; \Phi\bigl(\mathrm{IndCoh}_{\mathcal{N}}(\mathrm{LocSys}_{\LG}(C))\bigr)
\]
be the chiral algebra attached to ind-coherent sheaves on
$\LG$-local systems, with nilpotent singular support (Arinkin-Gaitsgory).
\end{enumerate}

\noindent\emph{Chiral Koszul conjecture.} There is a chiral Koszul
duality
\[
 \bigl(\mathcal{A}^{\mathrm{aut}}_{C,G}\bigr)^! \;=\; \mathcal{A}^{\mathrm{spec}}_{C,\LG}
\]
in $\Eone\text{-}\mathrm{ChirAlg}$, compatible with the Feigin-Frenkel
critical-level center isomorphism
\[
 Z\bigl(\widehat{\fg}_{\mathrm{crit}}\bigr) \;\cong\; \mathrm{Fun}(\Op_{\LG}(D)),
\]
where $\Op_{\LG}$ denotes $\LG$-opers on the formal disk $D$.

\noindent\emph{Claim status: CONJECTURAL} at CY scope (the right-hand
side of (b) is an $(\infty,1)$-categorical object whose chiralisation
is not yet constructed in full; morphism preservation for $\Phi$ is
AP-CY55-open). The $g = 0, 1$ cases reduce to Theorem \ref{thm:g0-yangian}
/ \ref{thm:g1-eht} and are consistent with Drinfeld-Kac-Moody
self-duality after rigidification.

\noindent\emph{Physical shadow (Kapustin-Witten).} The 4d $\mathcal{N}=4$
$S$-duality gives, on compactification on $\Sigma = \bS^2$:
$A$-branes $\leftrightarrow$ $B$-branes, i.e.\
$A$-branes on $\cM_H(C, G) \leftrightarrow B$-branes on
$\cM_H(C, \LG)$. The $\Phi$-chiralisation of this equivalence yields
Conjecture \ref{conj:chiral-langlands-koszul} modulo the (open) CY
chiralisation step.

\noindent\emph{Arithmetic shadow.} At $C$ a modular curve and
$G = \mathrm{GL}_2$, Conjecture \ref{conj:chiral-langlands-koszul}
intersects the Langlands reciprocity for $\mathrm{GL}_2(\mathbb{A}_{\bQ})$
only via the $p$-adic shadow tower (working notes line 2545,
HEURISTIC).
\end{conjecture}

\begin{remark}[On ``$G(C,G)^! = G(C, \LG)$'' as written in the manuscript]
The placeholder symbol $G(C,G)$ in working\_notes line 384 is the
``quantum vertex chiral group of the Hitchin system'', which is
precisely the object Conjecture \ref{conj:chiral-langlands-koszul}(a)
names as $\mathcal{A}^{\mathrm{aut}}_{C,G}$ after chiralisation;
Koszul-duality $(-)^!$ is the bar-cobar duality on
$\Eone$-chiral algebras. In this notation the slogan is correct, but
the manuscript should either (i) import Definition 3.4 of $G(C,G)$
from chapters/theory/quantum\_groups\_foundations.tex, or (ii) replace
the slogan with the explicit $\mathcal{A}^{\mathrm{aut}}/\mathcal{A}^{\mathrm{spec}}$
formulation above.
\end{remark}
\end{heal}

% ====================================================================
\section{Cluster B. Swiss-cheese and the $\Etwo$ bar complex (lines 1728--1738)}
% ====================================================================

\begin{attack}[A3: which Swiss-cheese operad?]\leavevmode

\noindent\emph{Ghost theorem.} The three bar complexes
$(B_{\Eone}, B_{\Etwo}, B_{\Einf})$ in working\_notes line 1731--1736
do correspond to three sectors of a two-coloured operad with open and
closed colours. Such an operad exists: Voronov's Swiss-cheese operad
$\mathcal{SC}_2$ (Voronov 1998, ``The Swiss-cheese operad'').

\noindent\emph{Precise error.} Two distinct operads share the name
``Swiss-cheese''. They are {\em not} interchangeable:

\begin{enumerate}[label=(\alph*)]
 \item \textbf{Voronov's $\mathcal{SC}_2$} (Voronov 1998) is a
   two-coloured topological operad whose closed colour is $E_2$ (little
   2-disks in the upper half-plane, moving in the interior) and whose
   open colour is $E_1$ (little intervals on the real boundary). An
   algebra over $\mathcal{SC}_2$ is a pair $(A, B)$ with $A$ an
   $E_2$-algebra, $B$ an $E_1$-algebra, and an $E_2$-action of $A$ on
   $B$ factoring through the Hochschild cohomology $A \to \mathrm{HH}^*(B,B)$.
   This is {\em Kontsevich's $(A, B)$-formality} / Deligne's conjecture
   setting.
 \item \textbf{Kontsevich's non-unital 2-disk operad} is the sub-operad
   of $E_2$ spanned by configurations with at least one disk. It is
   {\em not} two-coloured; it is single-colour $E_2$ minus the operadic
   unit. This is sometimes also called ``Swiss-cheese'' in physics
   lecture notes. It is {\em not} the operad controlling the
   open-closed sector structure.
 \item Costello-Gwilliam-Scheimbauer and the Vol II
   ``$\mathrm{SC}^{\mathrm{ch,top}}$'' specifically refer to a
   {\em chiral} refinement of Voronov's operad, carrying in addition
   a factorisation algebra structure in the Ran sense.
\end{enumerate}

The manuscript references ``Vol II Swiss-cheese structure
$\mathrm{SC}^{\mathrm{ch,top}}$ (Theorem H)'' (line 1731) without
importing the definition. A reader who follows the convention
``Swiss-cheese $=$ Voronov'' and another who follows ``Swiss-cheese
$=$ Kontsevich non-unital'' will read two different statements.

\noindent\emph{Correct relationship.} The operadic framework for
$(B_{\Eone}, B_{\Etwo}, B_{\Einf})$ is (c): the chiral Voronov
operad $\mathrm{SC}^{\mathrm{ch,top}}$, which is a two-coloured
$(\infty,1)$-operad in factorisation algebras with open colour $\Eone$
and closed colour $\Etwo$. The $\Einf$ bar complex is not part of
the Swiss-cheese structure per se; it lives at
$\Sym = \mathrm{colim}_{n\to\infty} E_n$, and the passage
$\Etwo \to \Einf$ is by the additional $\Einf$-delooping (Dunn), not
by the Swiss-cheese module action.
\end{attack}

\begin{heal}[H3: Voronov chiral refinement with labelled colours]\leavevmode

\begin{scope}[The operad controlling $(B_{\Eone}, B_{\Etwo}, B_{\Einf})$]\label{scope:swiss-cheese}
Let $\mathrm{SC}^{\mathrm{ch,top}}$ denote the two-coloured
$(\infty,1)$-factorisation operad (Costello-Gwilliam-Scheimbauer
version, formalised in Vol II Theorem H) with
\begin{itemize}
 \item closed colour: chiral $E_2$ (= $\Etwo$ with a compatible
   factorisation algebra on $C \in \Ran$);
 \item open colour: chiral $E_1$ (= $\Eone$ with compatible
   factorisation).
\end{itemize}
The structure maps are: operadic composition of closed disks
(closed sector $= \Etwo$); composition of open intervals on the
boundary (open sector $= \Eone$); an $\Etwo$-action of the closed
sector on the open sector (Deligne conjecture, chiral refinement).
\end{scope}

\begin{proposition}[Three bar complexes, three sectors; honest status]\label{prop:three-bar-sectors}
Let $A \in \Eone\text{-}\mathrm{ChirAlg}$ and let $(A, B_{\mathrm{op}})$
be a chiral $\mathrm{SC}^{\mathrm{ch,top}}$-algebra with closed sector
$B_{\mathrm{op}} \in \Etwo\text{-}\mathrm{ChirAlg}$ (equivalently,
$B_{\mathrm{op}} = \cZ^{\mathrm{der}}_{\mathrm{ch}}(A)$, Vol II
Theorem G). Then:
\begin{enumerate}[label=(\roman*)]
 \item $B_{\Eone}(A)$, the $\Eone$-bar complex of $A$, is the chiral
   Hochschild complex.
 \item $B_{\Etwo}(A) := B_{\Eone}(\cZ^{\mathrm{der}}_{\mathrm{ch}}(A))$,
   the $\Eone$-bar complex of the derived centre.
 \item $B_{\Einf}(A) := \mathrm{colim}_n B_{E_n}(A)$ exists when
   $A$ is $\Einf$-deloopable; not every $A$ admits this colimit.
\end{enumerate}
The $(B_{\Eone}, B_{\Etwo})$ pair is controlled by
$\mathrm{SC}^{\mathrm{ch,top}}$; the $B_{\Einf}$ layer is {\em not}
a Swiss-cheese sector (it lives in the symmetric / $\Einf$ limit),
and attributing it to the Swiss-cheese structure is an AP-CY-14
category-sliding.
\end{proposition}

\begin{remark}
Line 1737 of working\_notes asserts: ``The $\Eone \to \Etwo$ passage
via Drinfeld center (Theorem c3-drinfeld-center) corresponds to the
open-to-bulk passage in Vol II via the chiral derived center
$\cZ^{\mathrm{der}}_{\mathrm{ch}}(A)$.'' This is correct once
interpreted in Scope \ref{scope:swiss-cheese}: Drinfeld centre $=$
chiral derived centre at $d=3$ (Vol III Theorem C instantiation), and
the open-to-bulk passage in Voronov is the $\Eone$-sector $\to$
$\Etwo$-sector map sending $A \mapsto \cZ(A)$.  {\em Warning
against AP-CY57 narration}: the original sentence reads as
``corresponds to'', which is a synonym for ``we think these are
related''; replace by the explicit functor
$\mathrm{op\text{-}bulk}\colon \Eone\text{-ChirAlg} \to \Etwo\text{-ChirAlg}$,
$A \mapsto \cZ^{\mathrm{der}}_{\mathrm{ch}}(A)$.
\end{remark}
\end{heal}

% ====================================================================
\section{Cluster B continued. Completed modular Koszul datum (lines 1739--1748)}
% ====================================================================

\begin{attack}[A4: the completed modular Koszul datum for CY3]\leavevmode

\noindent\emph{Ghost theorem.} The eight-tuple
\[
 \Pi^{\mathrm{oc}}_X(A_{X_3}) =
 \bigl(\cF_X(A), \overline{B}_X(A), \Theta_A, L_A, (V^{\mathrm{br}}, T^{\mathrm{br}}),
 R^{\mathrm{mod}}_4, \cZ^{\mathrm{der}}_{\mathrm{ch}}(A), \Tr_A\bigr)
\]
has the structural shape of a {\em completed modular Koszul datum} in
the sense of Vol I: a chiral algebra $A$, its bar-cobar pair, an
$L_\infty$-twisting, a root lattice, a pair of modular forms
(Borcherds $V$ and theta $T$), a modular-group $R$-matrix, the derived
centre, and a trace. This is the right generic shape.

\noindent\emph{Precise error.} The statement (line 1742) that this
datum ``specializes to the CY3 setting'' {\em as eight components}
silently identifies $A_{X_3}$ with a specific family:

\begin{enumerate}[label=(\alph*)]
 \item $A_{X_3}$ unqualified is ambiguous: conifold, local $\bP^2$,
   quintic, $K3 \times E$, and the banana manifold are all CY$_3$'s
   with {\em different} Pi-oc data.
 \item ``BKM root system'' presupposes a $K3 \times E$ or a
   Monster-class CY$_3$; the quintic has no BKM root system in the
   sense relevant to $\Delta_5$.
 \item ``Borcherds denominator formula controls the bar Euler
   product'' is literally true at $K3 \times E$ (via Gritsenko's
   $\Delta_5$, working notes line 1746), but fails at the quintic
   (the topological string partition function $Z_{\mathrm{top}}(Q)$
   is {\em not} a Borcherds product; it is a Gopakumar-Vafa product,
   which is a different structure).
\end{enumerate}

\noindent\emph{Correct relationship.} $\Pi^{\mathrm{oc}}$ is a functor
from compact $\CY_3$ categories to eight-tuples, whose output shape is
fixed; the components $V^{\mathrm{br}}$ and $R^{\mathrm{mod}}_4$
depend on which CY$_3$. At $K3 \times E$, $V^{\mathrm{br}} = \Delta_5$
and the BKM root system is the Gritsenko-Nikulin-Scheithauer lattice
$II_{2,2} \oplus II_{0,1}$; at the quintic, $V^{\mathrm{br}}$ is the
MUM-monodromy conformal block and $R^{\mathrm{mod}}_4$ is the
Candelas-de la Ossa-Greene-Parkes $A$-model $R$-matrix. Stating the
datum generically requires the modifier ``with $X$-dependent
$V^{\mathrm{br}}, T^{\mathrm{br}}, R^{\mathrm{mod}}_4$''.
\end{attack}

\begin{heal}[H4: the $\Pi^{\mathrm{oc}}$ datum as a family-valued functor]\leavevmode

\begin{proposition}[$\Pi^{\mathrm{oc}}$ is family-valued]\label{prop:pi-oc-family}
There is an $(\infty,1)$-functor
\[
 \Pi^{\mathrm{oc}} \colon \mathrm{CY}\text{-}\mathrm{cat}_3
 \;\longrightarrow\;
 (\text{8-tuples, as in Vol~I Definition~X.Y}),
\]
$A_X \mapsto (\cF_X(A), \overline{B}_X(A), \Theta_A, L_A,
(V^{\mathrm{br}}_X, T^{\mathrm{br}}_X), R^{\mathrm{mod}}_{4,X},
\cZ^{\mathrm{der}}_{\mathrm{ch}}(A), \Tr_A)$,
where the last six components are {\em $X$-dependent}. Its values on
the canonical CY$_3$ landmarks:

\begin{center}
\renewcommand{\arraystretch}{1.2}
\small
\begin{tabular}{lcccc}
 \toprule
 $X$ & $L_A$ & $V^{\mathrm{br}}_X$ & $T^{\mathrm{br}}_X$ & status \\
 \midrule
 $K3 \times E$ & $II_{2,2} \oplus II_{0,1}$ & $\Delta_5$ & $\theta_{K3 \times E}$ & proved ({\em CY-A}$_{3}$) \\
 Conifold & $A_1$ & free-field log & $1$ & proved ({\em CY-A}$_{3}$) \\
 Local $\bP^2$ & $A_2$ & quantum toroidal $\fgl_1$ character & $\theta_{\bP^2}$ & proved \\
 Quintic & MUM & CDGP $A$-model & $\Theta_{\rm quintic}$ & conjectural (CY-B open) \\
 Banana ($\chi=0$) & trivial & uncurved & $1$ & proved (Costello-Li) \\
 \bottomrule
\end{tabular}
\end{center}
\noindent The ``BKM root system / Borcherds denominator formula
controls bar Euler product'' of line 1746 is the $K3 \times E$ row;
it does not generalise off that row.
\end{proposition}
\end{heal}

% ====================================================================
\section{Cluster C. Quiver-chart gluing, three answers (lines 2447--2493)}
% ====================================================================

\begin{attack}[A5: ``three answers'' are not all theorems in the same sense]\leavevmode

\noindent\emph{Ghost theorem.} The ``three answers'' for why $d=3$
gives $\Eone$ --- topological, structural, hexagon --- really do
correspond to three independent obstructions that vanish simultaneously
(topological) and are controlled by a common parameter (Serre duality,
odd vs.\ even $d$, Feigin-Frenkel signature) in a beautiful way.

\noindent\emph{Precise error.} The three items are {\em not} theorems
of the same status:

\begin{enumerate}[label=(\roman*)]
 \item ``$\pi_3(BU) = 0$'' is a \emph{theorem} (Bott periodicity, 1959).
 \item ``Antisymmetric Euler form prevents CoHA from carrying an
   $R$-matrix directly; $\Etwo$ braiding exists only on the Drinfeld
   centre'' is a \emph{theorem} (Schiffmann-Vasserot, Maulik-Okounkov;
   Vol III Theorem 5.2). The numerics
   ``$\mathcal{O}_2^{\mathrm{str}} = 0, 1, 27$'' are specific computed
   values.
 \item ``The deformation factor $D(\varepsilon_1, \varepsilon_2) =
   (\varepsilon_1 - \varepsilon_2)^2/(\varepsilon_1 + \varepsilon_2)^2$
   vanishes at self-dual, nonzero generic'' is a \emph{computation}
   (toric $\bC^3$, Nekrasov $\Omega$-background), which has not been
   shown to generalise to non-toric CY$_3$. It is a statement at toric
   scope, not a theorem at CY$_3$ scope. This is exactly AP-CY15:
   toric evidence $\not\Rightarrow$ compact CY$_3$ conclusion.
\end{enumerate}

Labelling all three as ``$\Eone \to \Etwo$ obstruction components''
with the same font is a scope-flattening (AP-CY2).

\noindent\emph{Correct relationship.} (i) is a universal topological
triviality; (ii) is a universal CY-categorical theorem (Serre duality
origin is genuinely universal); (iii) is toric-universal and
generalises {\em conjecturally} via the Maulik-Okounkov/K-theoretic
stable envelope to non-toric cases, with the {\em compact} CY$_3$
case open.

\end{attack}

\begin{heal}[H5: honest labelling of the three obstructions]\leavevmode

\begin{proposition}[Three obstructions, labelled scope]\label{prop:three-obstructions}
The obstruction to lifting an $\Eone$-chiral algebra $A_X$ attached to
a CY$_3$ category $D^b(X)$ to an $\Etwo$-chiral algebra decomposes as
\begin{enumerate}[label=\emph{(\roman*)}]
 \item Topological: $\pi_3(BU) = 0$ gives \emph{trivial} topological
   obstruction (Bott periodicity). Status: \textbf{universal theorem}.
 \item Structural: the antisymmetric Euler form for $d$ odd obstructs
   a direct CoHA $R$-matrix; $\Etwo$ braiding lives on
   $\cZ(\Rep^{\Eone}(A_X))$, not on $A_X$. Status: \textbf{universal theorem}
   (Schiffmann-Vasserot; Vol III \S5.2). First-principles origin:
   Serre duality $\chi(E,F) = (-1)^d\chi(F,E)$ + the identification
   $\pi_1(\mathrm{Conf}_2(\bR^{d-2}))$ trivial for $d-2$ odd, $\bZ$ for
   $d-2$ even.
 \item Hexagon-deformation: the toric deformation factor
   $D(\varepsilon_1,\varepsilon_2) = (\varepsilon_1 - \varepsilon_2)^2 /
   (\varepsilon_1 + \varepsilon_2)^2$ selects the self-dual locus. Status:
   \textbf{toric-scope theorem}, \emph{conjectural} at compact CY$_3$.
   The Maulik-Okounkov stable envelope provides the generalisation
   path; compact-case proof is open.
\end{enumerate}
The three are independent: (i) is topological, (ii) is categorical,
(iii) is parameter-space.
\end{proposition}

\begin{remark}[On the manuscript's ``Serre duality origin'' remark (line 2466)]
The manuscript's observation ``$\chi(E,F) = (-1)^d \chi(F,E)$ \ldots
algebraic shadow of $\pi_1(\mathrm{Conf}_2(\bR^{d-2}))$'' is correct
and sharp; this is the kind of sentence that earns its place. It
should be promoted to a named lemma (\emph{Serre-Fuks duality
coincidence lemma}) rather than remain a remark.
\end{remark}
\end{heal}

% ====================================================================
\section{Cluster C continued. CY$_4$ prediction and Dunn additivity (lines 2527--2534)}
% ====================================================================

\begin{attack}[A6: ``$\bS^4 = \bS^2 \times \bS^2$'' and Dunn additivity]\leavevmode

\noindent\emph{Ghost theorem.} The slogan $\CY_d \to E_{d-2}$ is a
genuine prediction: by Costello's $E_{d-2}$-scope theorem plus Dunn
additivity, a CY$_d$ category should give an $E_{d-2}$-chiral algebra
at the appropriate formality-preserved sense.

\noindent\emph{Precise error.} The line ``$\bS^4 = \bS^2 \times \bS^2$
(Hopf decomposition), each $\bS^2$ gives $\Eone$, two $\Eone$'s give
$\Etwo$ by Dunn additivity'' (line 2531) commits several sleights:

\begin{enumerate}[label=(\alph*)]
 \item $\bS^4 \not\cong \bS^2 \times \bS^2$ as manifolds:
   $\bS^2 \times \bS^2$ has $\chi = 4$, second Betti number $b_2 = 2$,
   and signature $0$; $\bS^4$ has $\chi = 2$, $b_2 = 0$, signature $0$.
   The correct statement is that $\bS^4$ is a {\em double-cover} of
   $\bS^2 \vee \bS^2$ via the Hopf \emph{suspension}, or more
   precisely, there is a smash-product decomposition
   $\bS^2 \wedge \bS^2 = \bS^4$. The manuscript's equals sign is
   wrong on the nose.
 \item ``Each $\bS^2$ gives $\Eone$'' is incorrect: $\bS^2$-framing
   gives $\Etwo$ (little $2$-disks), not $\Eone$.  The manuscript is
   thinking of the CY$_2 \to \Eone$ factor, but the $\bS^2$ in the
   Hopf decomposition refers to the framing sphere of the $E_n$ operad,
   not the CY dimension.
 \item ``Two $\Eone$'s give $\Etwo$ by Dunn additivity'' is correct as
   a statement ($\Eone \otimes \Eone = \Etwo$; Dunn 1988, Lurie HA
   Theorem 5.1.2.2), but not applicable here: the two $\Eone$'s
   would need to come from two commuting $\Eone$-structures on the
   same algebra, not from a product decomposition of the framing sphere.
\end{enumerate}

\noindent\emph{Correct relationship.} The CY$_4 \to \Etwo$ prediction
should be stated as: for $X$ a compact smooth CY$_4$, $\Phi(A_X)$ is an
$\Etwo$-chiral algebra by the $E_{d-2}$-output-scope rule ($d = 4$
gives $E_2$, i.e.\ $\Etwo$). Dunn additivity enters the proof path
as: $\Etwo = \Eone \otimes \Eone$ on the level of operads, and the two
commuting $\Eone$-structures on $\Phi(A_X)$ come from (i) the
$C^\infty$ symplectic form on $X$ and (ii) the holomorphic volume form
$\Omega_X$ (complex 4-form). These are compatible because $X$ is
$K$-trivial and hyperkähler-like in the generalised sense.
\end{attack}

\begin{heal}[H6: CY$_4 \to \Etwo$ via Dunn additivity of two commuting $\Eone$-structures]\leavevmode

\begin{conjecture}[CY$_4$ gives $\Etwo$]\label{conj:cy4-e2}
Let $X$ be a compact smooth $K$-trivial complex 4-fold (CY$_4$) with
$K_X \cong \cO_X$. Then the CY-to-chiral functor $\Phi$ sends
$D^b(\Coh(X))$ to an $\Etwo$-chiral algebra
\[
 \Phi \colon D^b(\Coh(X)) \;\longmapsto\; \mathcal{A}^{(4)}_X
 \;\in\; \Etwo\text{-}\mathrm{ChirAlg}.
\]
The $\Etwo$ structure is witnessed by Dunn additivity
$\Etwo = \Eone \otimes \Eone$ applied to the pair of commuting
$\Eone$-structures on $\mathcal{A}^{(4)}_X$ coming from:
\begin{itemize}
 \item the $C^\infty$ symplectic $\Eone$ (arising from the
   Nakajima-type raising-lowering operators in CoHA);
 \item the holomorphic $\Eone$ (arising from the 4-form
   $\Omega_X$ via Costello-Gwilliam factorisation).
\end{itemize}
The two $\Eone$'s commute by the degeneracy of the
Fröhlicher-type spectral sequence on $X$ (Hodge-symmetric type).

\noindent\emph{Claim status: CONJECTURAL.} Supporting cases: $X = K3 \times K3$
(where the two $\Eone$'s come from the two K3 factors; proof
{\em reduces} to chiral Koszul on $K3$, which is Vol III \S7), and
$X$ hyperkähler (where a single hyperkähler $\bS^2$-family of complex
structures gives an $\bS^2$-framed $\Etwo$ directly, Kapustin-Rozansky-Saulina
style). The general smooth CY$_4$ case is open.
\end{conjecture}

\begin{remark}[Correction on $\bS^4 \neq \bS^2 \times \bS^2$]
The working-notes line 2531 should be corrected to: ``For
$X = K3 \times K3$: $X$ has two commuting $\Eone$-structures, one per
$K3$ factor. By Dunn additivity $\Eone \otimes \Eone = \Etwo$, the
product carries an $\Etwo$-structure.'' The sphere-framing sentence
about $\bS^4$ vs.~$\bS^2 \times \bS^2$ is not needed for the
argument; it conflates the framing sphere of the $E_n$-operad with
the manifold $X$. (AP-CY14 category-sliding.)
\end{remark}
\end{heal}

% ====================================================================
\section{Cluster C final. Quantum toroidal from 3-chart hocolim (lines 2551--2566)}
% ====================================================================

\begin{attack}[A7: hocolim vs Hall product]\leavevmode

\noindent\emph{Ghost theorem.} For local $\bP^2$ the three Seiberg-dual
quiver charts $(Q_0, Q_1, Q_2)$ do carry a $\bZ_3$-action compatible
with the Miki automorphism of the quantum toroidal
$U_{q,t}(\widehat{\widehat{\fgl}}_1)$, and an $\Eone$-chiral hocolim in
$\Eone$-ChirAlg does compute the quantum toroidal algebra at the level
stated (through total degree 6).

\noindent\emph{Precise error.} The statement (line 2551--2566)
``$U_{q,t}(\widehat{\widehat{\fgl}}_1) = \mathrm{hocolim}(\CoHA(Q_0,W_0)
\rightrightarrows \CoHA(Q_1,W_1) \rightrightarrows \CoHA(Q_2,W_2))$''
has three subtle issues:

\begin{enumerate}[label=(\alph*)]
 \item The hocolim is in \emph{which} $(\infty,1)$-category? Plain
   $\Eone$-algebras? $\Eone$-ChirAlg? The Ran-factorisation category?
   Each gives a different answer.
 \item The equality is through total degree $6$ only, per the
   manuscript, but the symbol ``$=$'' is used. This should be
   $\simeq_{\leq 6}$ (equivalence after truncating to total degree
   $\leq 6$).
 \item Sequential mutations do not compose (exchange matrix
   entries grow exponentially); the manuscript acknowledges this
   (line 2551), but then does not explain what operadic diagram
   actually drives the hocolim. The answer is: the cyclic diagram
   $Q_0 \to Q_1 \to Q_2 \to Q_0$ in the $(\infty,1)$-category of
   $\Eone$-ChirAlg{\em's} with conjugation-mutation morphisms, not
   sequential mutations.
\end{enumerate}

\noindent\emph{Correct relationship.} The correct statement is an
equivalence of $\Eone$-chiral algebras up to truncation:
\[
 \tau_{\leq 6}\bigl(U_{q,t}(\widehat{\widehat{\fgl}}_1)\bigr)
 \;\simeq\;
 \tau_{\leq 6}\bigl(\mathrm{hocolim}^{\Eone\text{-ChirAlg}}
 (\CoHA(Q_0,W_0) \rightrightarrows \CoHA(Q_1,W_1)
 \rightrightarrows \CoHA(Q_2,W_2))\bigr),
\]
where the two parallel arrows are the two conjugation-mutation
morphisms (not sequential compositions), the cyclic diagram has three
objects forming a $\bZ_3$-orbit, and the Miki automorphism $S$
satisfies $S^3 = \mathrm{id}$ on the hocolim.
\end{attack}

\begin{heal}[H7: truncated hocolim in $\Eone$-ChirAlg, $\bZ_3$-Miki]\leavevmode

\begin{conjecture}[Quantum toroidal as 3-chart hocolim, truncated]\label{conj:qt-hocolim}
Let local $\bP^2$ be the CY$_3$ with quiver description given by the
three Seiberg-dual quivers $Q_0, Q_1, Q_2$ obtained by mutating the
original $Q_0$ at its three nodes; let $W_i$ denote the superpotential
transported to $Q_i$ by conjugation, not by sequential mutation. Then
there is an $\Eone$-chiral equivalence up to degree-6 truncation
\[
 \tau_{\leq 6}\,U_{q,t}\bigl(\widehat{\widehat{\fgl}}_1\bigr)
 \;\simeq\;
 \tau_{\leq 6}\,\mathrm{hocolim}^{\Eone\text{-ChirAlg}}_{\bZ_3}
 \bigl(\CoHA(Q_0,W_0)
 \rightrightarrows \CoHA(Q_1,W_1)
 \rightrightarrows \CoHA(Q_2,W_2)\bigr),
\]
with the $\bZ_3$-action of the hocolim descending to the Miki
automorphism $S \in \mathrm{Aut}(U_{q,t})$ satisfying $S^3 = \mathrm{id}$.
The bar-hocolim commutation
\[
 B^{\Eone}(U_{q,t}) \simeq \mathrm{hocolim}\, B^{\Eone}(\CoHA_i)
\]
holds at the same truncation level.

\emph{Claim status: CONJECTURAL at arbitrary total degree}; proved at
total degree $\leq 6$ by the 124-test engine
\texttt{quantum\_toroidal\_from\_3chart.py}. Extending beyond degree
6 is open. The mechanism is similar to Maulik-Okounkov stable envelope
$\bZ_3$-invariants for equivariant K-theory of the resolved conifold;
the extension to non-toric compact CY$_3$ remains the object of
Pattern 236 ambient-qualifier discipline.
\end{conjecture}
\end{heal}

% ====================================================================
\section*{Cache: new AP-CYs proposed}
% ====================================================================

The investigation surfaces three new antipatterns for inclusion in
\texttt{notes/antipatterns\_catalogue.md}:

\begin{itemize}
 \item \textbf{AP-CY63 --- Langlands scope conflation (Critical).}
   ``Langlands'' is tri-valued at CY scope: arithmetic (Gal $\leftrightarrow$
   automorphic), Beilinson-Drinfeld geometric ($\DMod(\Bun_G) \simeq
   \mathrm{IndCoh}_\mathcal{N}(\mathrm{LocSys}_{\LG})$), Kapustin-Witten physical
   (4d $\mathcal{N}=4$ $S$-duality). Each use of ``Langlands'' must
   name which. \emph{Counter}: the CY-scope operative statement is
   usually (L2), with (L3) as physical shadow and (L1) as $p$-adic
   consequence; the manuscript conflates all three in the slogan
   $G(C,G)^! = G(C,\LG)$.

 \item \textbf{AP-CY64 --- Swiss-cheese operad ambiguity (High).}
   ``Swiss-cheese'' refers to Voronov's two-coloured 2-disk operad
   (open colour $\Eone$, closed colour $\Etwo$), or Kontsevich's
   non-unital single-colour 2-disk operad, or the chiral refinement
   $\mathrm{SC}^{\mathrm{ch,top}}$ of Vol II. These are three different
   operads. \emph{Counter}: always name the reference (Voronov 1998 /
   Kontsevich 1999 / Vol II Theorem H).

 \item \textbf{AP-CY65 --- $\bS^4 = \bS^2 \times \bS^2$ is false
   (Low-Medium).} The Hopf decomposition gives $\bS^2 \wedge \bS^2 = \bS^4$,
   not a Cartesian product. When using Dunn additivity for CY$_4 \to \Etwo$,
   the two commuting $\Eone$-structures come from the pair
   (symplectic, holomorphic), not from a product decomposition of the
   framing sphere. \emph{Counter}: state Dunn additivity in operadic
   form $\Eone \otimes \Eone = \Etwo$, not as a topological
   decomposition of framing data.
\end{itemize}

% ====================================================================
\section*{Summary of repairs}
% ====================================================================

\begin{itemize}
 \item Working-notes \textbf{lines 364--375}: the ``$\CY_2$ escape'' is
   a functor on the Hitchin moduli space $\cM_H(C,G)$, not on
   $\Coh(T^*C)$ literally; the slogan
   ``$\CY_3 = \CY_2 + \text{automorphic correction}$'' is a
   denominator-formula ($\kBKM$) statement at $K3 \times E$, not a
   $\kch$ / $\kcat$ / $\kappa_{\mathrm{fiber}}$ invariant statement.
 \item Working-notes \textbf{lines 378--389} (``Langlands''): the
   conjecture is (L2) Beilinson-Drinfeld, stated as chiral Koszul
   duality $(\mathcal{A}^{\mathrm{aut}})^! = \mathcal{A}^{\mathrm{spec}}$; (L1)
   arithmetic and (L3) Kapustin-Witten are shadows.
 \item Working-notes \textbf{line 1731--1737}: ``Swiss-cheese'' means
   chiral Voronov $\mathrm{SC}^{\mathrm{ch,top}}$; $B_{\Einf}$ is not
   a Swiss-cheese sector; open-to-bulk passage is the explicit
   functor $\cZ^{\mathrm{der}}_{\mathrm{ch}}$.
 \item Working-notes \textbf{line 1742--1746}: $\Pi^{\mathrm{oc}}$
   is family-valued; ``BKM root system / Borcherds denominator'' is
   the $K3 \times E$ row, not universal over CY$_3$.
 \item Working-notes \textbf{line 2459--2466}: three
   $\Eone \to \Etwo$ obstructions are not all universal theorems; (i)
   topological is universal, (ii) structural is universal, (iii)
   hexagon-deformation is toric-scope with compact generalisation
   open.
 \item Working-notes \textbf{line 2531}: $\bS^4 \neq \bS^2 \times \bS^2$;
   the CY$_4 \to \Etwo$ argument uses Dunn additivity on two
   commuting $\Eone$-structures (symplectic + holomorphic), not a
   product decomposition.
 \item Working-notes \textbf{line 2551--2566}: equality of hocolim
   with $U_{q,t}$ holds at truncation $\leq 6$; the cyclic diagram
   uses conjugation-mutation, not sequential mutation; $\bZ_3$-action
   is the Miki automorphism.
\end{itemize}

\end{document}
